% Prajñāpāramitāhṛdaya - Heart Sūtra
% Portrait format with per-page footnotes
% Three columns, phrase-level units

\documentclass[9pt,a4paper]{article}

% ============================================================================
% PACKAGES
% ============================================================================

\usepackage[margin=1.5cm,bottom=2.5cm]{geometry}
\usepackage{fontspec}
\usepackage{xeCJK}
\usepackage{tabularx}
\usepackage{array}
\usepackage{booktabs}
\usepackage{xcolor}
\usepackage{longtable}

% Per-page footnotes
\usepackage[perpage,symbol*]{footmisc}
\renewcommand{\thefootnote}{\arabic{footnote}}

% Page style
\usepackage{fancyhdr}
\pagestyle{fancy}
\fancyhf{}
\renewcommand{\headrulewidth}{0pt}
\fancyfoot[C]{\thepage}

% ============================================================================
% FONTS
% ============================================================================

\setmainfont{Noto Serif}
\setsansfont{Noto Sans}
\setCJKmainfont{Noto Serif CJK SC}
\setCJKsansfont{Noto Sans CJK SC}

\newfontfamily\devanagarifont{Noto Serif Devanagari}
\newcommand{\deva}[1]{{\devanagarifont #1}}

\newfontfamily\tibetanfont{Noto Serif Tibetan}
% Tibetan smaller with flexible breaks between syllables
\newcommand{\tib}[1]{{\tibetanfont\scriptsize\tolerance=9999\emergencystretch=3em #1}}

% ============================================================================
% CUSTOM COMMANDS
% ============================================================================

% Line number
\newcommand{\linenum}[1]{\textbf{\small #1}}

% Romanization (smaller, gray)
\newcommand{\rom}[1]{{\scriptsize\color{gray} #1}}

% English translation - spans all columns, CENTERED, with space after
\newcommand{\eng}[1]{\multicolumn{3}{@{}c@{}}{\footnotesize\itshape #1} \\[0.5em]}

% Apparatus marker - uses footnotes
\newcommand{\var}[1]{\textsuperscript{\color{red}\textbf{#1}}}

% For apparatus
\newcommand{\lemma}[1]{\textbf{#1}}
\newcommand{\reading}[1]{#1}
\newcommand{\wit}[1]{\textsc{#1}}

% ============================================================================
% DOCUMENT
% ============================================================================

\begin{document}

% ----------------------------------------------------------------------------
% TITLE
% ----------------------------------------------------------------------------

\begin{center}
{\Large\bfseries Prajñāpāramitāhṛdaya}~~般若波羅蜜多心經~~\deva{प्रज्ञापारमिताहृदय}\\[0.3em]
{\large Heart Sūtra -- Synoptic Critical Edition}\\[0.2em]
{\small Chinese T251 (649 CE) · Sanskrit GRETIL · Tibetan Toh 21}
\end{center}

\vspace{0.5em}
\hrule
\vspace{0.5em}

% ----------------------------------------------------------------------------
% TEXT
% ----------------------------------------------------------------------------

\footnotesize

% Column widths: Chinese 26%, Sanskrit 32%, Tibetan 38% (Tibetan is verbose)
\begin{longtable}{@{}p{0.26\textwidth}|p{0.32\textwidth}|p{0.38\textwidth}@{}}

\textbf{Chinese} & \textbf{Sanskrit} & \textbf{Tibetan} \\
\hline
\endfirsthead
\textbf{Chinese} & \textbf{Sanskrit} & \textbf{Tibetan} \\
\hline
\endhead

% === LINE 1 ===
\linenum{1} 觀自在\footnote{\lemma{觀自在}] \reading{觀世音} \wit{T250 T252}. Supports Skt. \textit{-īśvara} ``lord'' vs. \textit{-svara} ``sound.''}菩薩
\newline\rom{Guānzìzài púsà}
&
\deva{आर्यावलोकितेश्वरो}\footnote{\lemma{avalokiteśvara-}] \reading{avalokita-} \wit{Ja}. Chinese 觀自在 supports \textit{-īśvara}.}\deva{बोधिसत्त्वो}
\newline\rom{Āryāvalokiteśvaro bodhisattvo}
&
\tib{འཕགས་པ་སྤྱན་རས་གཟིགས་དབང་ཕྱུག་བྱང་ཆུབ་སེམས་དཔའ}
\newline\rom{'phags pa spyan ras gzigs dbang phyug byang chub sems dpa'}
\\
\eng{The noble Avalokiteśvara Bodhisattva}

% === LINE 2 ===
\linenum{2} 行深般若波羅蜜多時
\newline\rom{xíng shēn bōrě bōluómìduō shí}
&
\deva{गम्भीरां प्रज्ञापारमिताचर्यां चरमाणो}
\newline\rom{gambhīrāṃ prajñāpāramitā-caryāṃ caramāṇo}
&
\tib{ཤེས་རབ་ཀྱི་ཕ་རོལ་ཏུ་ཕྱིན་པ་ཟབ་མོའི་སྤྱོད་པ་སྤྱད་པ}
\newline\rom{shes rab kyi pha rol tu phyin pa zab mo'i spyod pa spyad pa}
\\
\eng{practicing the profound prajñāpāramitā}

% === LINE 3 ===
\linenum{3} 照見五蘊皆空
\newline\rom{zhàojiàn wǔyùn jiē kōng}
&
\deva{व्यवलोकयति स्म}\footnote{\lemma{vyavalokayati sma}] Unusual imperfect. May reflect Chinese narrative 時 \textit{shí}.} \deva{पञ्चस्कन्धास् तांश्च स्वभावशून्यान् पश्यति स्म}
\newline\rom{vyavalokayati sma pañca-skandhās tāṃś ca svabhāva-śūnyān paśyati sma}
&
\tib{ཕུང་པོ་ལྔ་པོ་དེ་དག་ལ་ཡང་རང་བཞིན་གྱིས་སྟོང་པར་རྣམ་པར་བལྟའོ}
\newline\rom{phung po lnga po de dag la yang rang bzhin gyis stong par rnam par blta'o}
\\
\eng{illuminated and saw that the five aggregates are all empty}

% === LINE 4 ===
\linenum{4} 度一切苦厄\footnote{\lemma{度一切苦厄}] No Sanskrit parallel. Chinese uniquely adds ``crossed beyond all suffering.''}
\newline\rom{dù yīqiè kǔ è}
&
---
\newline\rom{[not in Sanskrit]}
&
---
\newline\rom{[not in short version]}
\\
\eng{and crossed beyond all suffering}
\hline

% === LINE 5 ===
\linenum{5} 舍利子
\newline\rom{Shèlìzǐ}
&
\deva{इह शारिपुत्र}
\newline\rom{Iha Śāriputra}
&
\tib{ཤཱ་རིའི་བུ}
\newline\rom{shA ri'i bu}
\\
\eng{Śāriputra,}

% === LINE 6 ===
\linenum{6} 色不異空
\newline\rom{sè bù yì kōng}
&
\deva{रूपं शून्यता}
\newline\rom{rūpaṃ śūnyatā}
&
\tib{གཟུགས་སྟོང་པའོ}
\newline\rom{gzugs stong pa'o}
\\
\eng{form is not different from emptiness,}

% === LINE 7 ===
\linenum{7} 空不異色
\newline\rom{kōng bù yì sè}
&
\deva{शून्यतैव रूपम्}
\newline\rom{śūnyataiva rūpam}
&
\tib{སྟོང་པ་ཉིད་གཟུགས་སོ}
\newline\rom{stong pa nyid gzugs so}
\\
\eng{emptiness is not different from form.}

% === LINE 8 ===
\linenum{8} 色即是空，空即是色
\newline\rom{sè jí shì kōng, kōng jí shì sè}
&
\deva{रूपान् न पृथक् शून्यता शून्यतायाः न पृथग् रूपम्}
\newline\rom{rūpān na pṛthak śūnyatā śūnyatāyā na pṛthag rūpam}
&
\tib{གཟུགས་ལས་སྟོང་པ་ཉིད་གཞན་མ་ཡིན། སྟོང་པ་ཉིད་ལས་ཀྱང་གཟུགས་གཞན་མ་ཡིན་ནོ}
\newline\rom{gzugs las stong pa nyid gzhan ma yin. stong pa nyid las kyang gzugs gzhan ma yin no}
\\
\eng{Form is emptiness, emptiness is form.}

% === LINE 9 ===
\linenum{9} 受想行識亦復如是
\newline\rom{shòu xiǎng xíng shí yì fù rúshì}
&
\deva{एवमेव वेदनासंज्ञासंस्कारविज्ञानानि}
\newline\rom{evam eva vedanā-saṃjñā-saṃskāra-vijñānāni}
&
\tib{དེ་བཞིན་དུ་ཚོར་བ་དང་འདུ་ཤེས་དང་འདུ་བྱེད་དང་རྣམ་པར་ཤེས་པ་རྣམས་སྟོང་པའོ}
\newline\rom{de bzhin du tshor ba dang 'du shes dang 'du byed dang rnam par shes pa rnams stong pa'o}
\\
\eng{Feeling, perception, volition, and consciousness are also like this.}
\hline

% === LINE 10 ===
\linenum{10} 舍利子，是諸法空相
\newline\rom{Shèlìzǐ, shì zhūfǎ kōng xiāng}
&
\deva{इह शारिपुत्र सर्वधर्माः शून्यतालक्षणा}
\newline\rom{Iha Śāriputra sarva-dharmāḥ śūnyatā-lakṣaṇā}
&
\tib{ཤཱ་རིའི་བུ། དེ་ལྟར་ཆོས་ཐམས་ཅད་སྟོང་པ་ཉིད་དེ། མཚན་ཉིད་མེད་པ}
\newline\rom{shA ri'i bu. de ltar chos thams cad stong pa nyid de. mtshan nyid med pa}
\\
\eng{Śāriputra, all dharmas are marked by emptiness:}

% === LINE 11 ===
\linenum{11} 不生不滅
\newline\rom{bù shēng bù miè}
&
\deva{अनुत्पन्ना अनिरुद्धा}
\newline\rom{anutpannā aniruddhā}
&
\tib{མ་སྐྱེས་པ། མ་འགགས་པ}
\newline\rom{ma skyes pa. ma 'gags pa}
\\
\eng{not arising, not ceasing;}

% === LINE 12 ===
\linenum{12} 不垢不淨
\newline\rom{bù gòu bù jìng}
&
\deva{अमला अविमला}
\newline\rom{amalā avimalā}
&
\tib{དྲི་མ་མེད་པ། དྲི་མ་དང་བྲལ་བ་མེད་པ}
\newline\rom{dri ma med pa. dri ma dang bral ba med pa}
\\
\eng{not defiled, not pure;}

% === LINE 13 ===
\linenum{13} 不增不減
\newline\rom{bù zēng bù jiǎn}
&
\deva{अनूना अपरिपूर्णाः}
\newline\rom{anūnā aparipūrṇāḥ}
&
\tib{བྲི་བ་མེད་པ། གང་བ་མེད་པའོ}
\newline\rom{bri ba med pa. gang ba med pa'o}
\\
\eng{not increasing, not decreasing.}
\hline

% === LINE 14 ===
\linenum{14} 是故空中，無色
\newline\rom{Shìgù kōng zhōng, wú sè}
&
\deva{तस्माच्छारिपुत्र शून्यतायां न रूपं}
\newline\rom{Tasmāc Chāriputra śūnyatāyāṃ na rūpaṃ}
&
\tib{དེ་ལྟ་བས་ན་སྟོང་པ་ཉིད་ལ་གཟུགས་མེད}
\newline\rom{de lta bas na stong pa nyid la gzugs med}
\\
\eng{Therefore in emptiness there is no form,}

% === LINE 15 ===
\linenum{15} 無受想行識
\newline\rom{wú shòu xiǎng xíng shí}
&
\deva{न वेदना न संज्ञा न संस्काराः न विज्ञानम्}
\newline\rom{na vedanā na saṃjñā na saṃskārāḥ na vijñānam}
&
\tib{ཚོར་བ་མེད། འདུ་ཤེས་མེད། འདུ་བྱེད་རྣམས་མེད། རྣམ་པར་ཤེས་པ་མེད}
\newline\rom{tshor ba med. 'du shes med. 'du byed rnams med. rnam par shes pa med}
\\
\eng{no feeling, perception, volition, or consciousness;}

% === LINE 16 ===
\linenum{16} 無眼耳鼻舌身意
\newline\rom{wú yǎn ěr bí shé shēn yì}
&
\deva{न चक्षुःश्रोत्रघ्राणजिह्वाकायमनांसि}
\newline\rom{na cakṣuḥ-śrotra-ghrāṇa-jihvā-kāya-manāṃsi}
&
\tib{མིག་མེད། རྣ་བ་མེད། སྣ་མེད། ལྕེ་མེད། ལུས་མེད། ཡིད་མེད}
\newline\rom{mig med. rna ba med. sna med. lce med. lus med. yid med}
\\
\eng{no eye, ear, nose, tongue, body, or mind;}

% === LINE 17 ===
\linenum{17} 無色聲香味觸法
\newline\rom{wú sè shēng xiāng wèi chù fǎ}
&
\deva{न रूपशब्दगन्धरसस्प्रष्टव्यधर्माः}
\newline\rom{na rūpa-śabda-gandha-rasa-spraṣṭavya-dharmāḥ}
&
\tib{གཟུགས་མེད། སྒྲ་མེད། དྲི་མེད། རོ་མེད། རེག་བྱ་མེད། ཆོས་མེད་དོ}
\newline\rom{gzugs med. sgra med. dri med. ro med. reg bya med. chos med do}
\\
\eng{no form, sound, smell, taste, touch, or mental objects;}

% === LINE 18 ===
\linenum{18} 無眼界，乃至無意識界
\newline\rom{wú yǎn jiè, nǎizhì wú yìshí jiè}
&
\deva{न चक्षुर्धातुर् यावन् न मनोविज्ञानधातुः}
\newline\rom{na cakṣur-dhātur yāvan na manovijñāna-dhātuḥ}
&
\tib{མིག་གི་ཁམས་མེད་པ་ནས་ཡིད་ཀྱི་རྣམ་པར་ཤེས་པའི་ཁམས་ཀྱི་བར་དུ་ཡང་མེད་དོ}
\newline\rom{mig gi khams med pa nas yid kyi rnam par shes pa'i khams kyi bar du yang med do}
\\
\eng{no eye-element up to no mind-consciousness element.}
\hline

% === LINE 19 ===
\linenum{19} 無無明，亦無無明盡\footnote{\lemma{盡} \textit{jìn}] = Skt. \textit{kṣaya} ``destruction,'' not standard \textit{nirodha} ``cessation.'' Key back-translation evidence.}
\newline\rom{wú wúmíng, yì wú wúmíng jìn}
&
\deva{न अविद्या न अविद्याक्षयो}\footnote{\lemma{avidyā-kṣayo}] Uses \textit{kṣaya} where Large PP has \textit{nirodha}. ``Smoking gun'' evidence (Nattier 1992:175--8).}
\newline\rom{na-avidyā na-avidyā-kṣayo}
&
\tib{མ་རིག་པ་མེད། མ་རིག་པ་ཟད་པ་མེད་པ་ནས}
\newline\rom{ma rig pa med. ma rig pa zad pa med pa nas}
\\
\eng{No ignorance and no ending of ignorance,}

% === LINE 20 ===
\linenum{20} 乃至無老死，亦無老死盡
\newline\rom{nǎizhì wú lǎosǐ, yì wú lǎosǐ jìn}
&
\deva{यावन् न जरामरणं न जरामरणक्षयः}
\newline\rom{yāvan na jarā-maraṇaṃ na jarā-maraṇa-kṣayaḥ}
&
\tib{རྒ་ཤི་མེད། རྒ་ཤི་ཟད་པའི་བར་དུ་ཡང་མེད་དོ}
\newline\rom{rga shi med. rga shi zad pa'i bar du yang med do}
\\
\eng{up to no old age and death, and no ending of old age and death.}

% === LINE 21 ===
\linenum{21} 無苦集滅道
\newline\rom{wú kǔ jí miè dào}
&
\deva{न दुःखसमुदयनिरोधमार्गाः}
\newline\rom{na duḥkha-samudaya-nirodha-mārgāḥ}
&
\tib{སྡུག་བསྔལ་བ་དང་། ཀུན་འབྱུང་བ་དང་། འགོག་པ་དང་། ལམ་མེད}
\newline\rom{sdug bsngal ba dang. kun 'byung ba dang. 'gog pa dang. lam med}
\\
\eng{No suffering, origination, cessation, or path.}

% === LINE 22 ===
\linenum{22} 無智亦無得，以無所得故
\newline\rom{wú zhì yì wú dé, yǐ wú suǒ dé gù}
&
\deva{न ज्ञानं न प्राप्तिर् न अप्राप्तिः}
\newline\rom{na jñānaṃ na prāptir na-aprāptiḥ}
&
\tib{ཡེ་ཤེས་མེད། ཐོབ་པ་མེད། མ་ཐོབ་པ་ཡང་མེད་དོ}
\newline\rom{ye shes med. thob pa med. ma thob pa yang med do}
\\
\eng{No wisdom and no attainment---because there is nothing to attain.}
\hline

% === LINE 23 ===
\linenum{23} 菩提薩埵，依般若波羅蜜多故
\newline\rom{Pútísàduǒ, yī bōrě bōluómìduō gù}
&
\deva{तस्माच्छारिपुत्र अप्राप्तित्वाद् बोधिसत्त्वस्य प्रज्ञापारमितामाश्रित्य विहरत्यचित्तावरणः}
\newline\rom{Tasmāc Chāriputra aprāptitvād bodhisattvasya prajñāpāramitām āśritya viharaty acittāvaraṇaḥ}
&
\tib{བྱང་ཆུབ་སེམས་དཔའ་རྣམས་ཐོབ་པ་མེད་པའི་ཕྱིར། ཤེས་རབ་ཀྱི་ཕ་རོལ་ཏུ་ཕྱིན་པ་ལ་བརྟེན་ཅིང་གནས་ཏེ}
\newline\rom{byang chub sems dpa' rnams thob pa med pa'i phyir. shes rab kyi pha rol tu phyin pa la brten cing gnas te}
\\
\eng{Bodhisattvas, relying on prajñāpāramitā,}

% === LINE 24 ===
\linenum{24} 心無罣礙，無罣礙故
\newline\rom{xīn wú guà'ài, wú guà'ài gù}
&
\deva{चित्तावरणनास्तित्वाद्}
\newline\rom{cittāvaraṇa-nāstitvād}
&
\tib{སེམས་ལ་སྒྲིབ་པ་མེད་པས}
\newline\rom{sems la sgrib pa med pas}
\\
\eng{have minds without obstruction.}

% === LINE 25 ===
\linenum{25} 無有恐怖，遠離顚倒夢想
\newline\rom{wú yǒu kǒngbù, yuǎnlí diāndǎo mèngxiǎng}
&
\deva{अत्रस्त्रो विपर्यासातिक्रान्तो}
\newline\rom{atrastro viparyāsa-atikrānto}
&
\tib{འཇིགས་པ་མེད་དེ། ཕྱིན་ཅི་ལོག་ལས་ཤིན་ཏུ་འདས་ནས}
\newline\rom{'jigs pa med de. phyin ci log las shin tu 'das nas}
\\
\eng{Without obstruction, they have no fear, are far from inverted dreams,}

% === LINE 26 ===
\linenum{26} 究竟涅槃
\newline\rom{jiūjìng nièpán}
&
\deva{निष्ठानिर्वाणप्राप्तः}\footnote{\lemma{niṣṭhā-nirvāṇa}] Controversial compound (Attwood 2018). Chinese 究竟涅槃 is clear.}
\newline\rom{niṣṭhā-nirvāṇa-prāptaḥ}
&
\tib{མྱ་ངན་ལས་འདས་པའི་མཐར་ཕྱིན་ཏོ}
\newline\rom{mya ngan las 'das pa'i mthar phyin to}
\\
\eng{and attain ultimate nirvāṇa.}
\hline

% === LINE 27 ===
\linenum{27} 三世諸佛，依般若波羅蜜多故
\newline\rom{Sānshì zhū fó, yī bōrě bōluómìduō gù}
&
\deva{त्र्यध्वव्यवस्थिताः सर्वबुद्धाः प्रज्ञापारमितामाश्रित्य}
\newline\rom{tryadhva-vyavasthitāḥ sarva-buddhāḥ prajñāpāramitām āśritya}
&
\tib{དུས་གསུམ་དུ་རྣམ་པར་བཞུགས་པའི་སངས་རྒྱས་ཐམས་ཅད་ཀྱང་ཤེས་རབ་ཀྱི་ཕ་རོལ་ཏུ་ཕྱིན་པ་ལ་བརྟེན་ནས}
\newline\rom{dus gsum du rnam par bzhugs pa'i sangs rgyas thams cad kyang shes rab kyi pha rol tu phyin pa la brten nas}
\\
\eng{All buddhas of the three times, relying on prajñāpāramitā,}

% === LINE 28 ===
\linenum{28} 得阿耨多羅三藐三菩提
\newline\rom{dé ānòuduōluó sānmiǎo sānpútí}
&
\deva{अनुत्तरां सम्यक्संबोधिमभिसंबुद्धाः}
\newline\rom{anuttarāṃ samyak-saṃbodhim abhisaṃbuddhāḥ}
&
\tib{བླ་ན་མེད་པ་ཡང་དག་པར་རྫོགས་པའི་བྱང་ཆུབ་ཏུ་མངོན་པར་རྫོགས་པར་སངས་རྒྱས་སོ}
\newline\rom{bla na med pa yang dag par rdzogs pa'i byang chub tu mngon par rdzogs par sangs rgyas so}
\\
\eng{attain anuttarā-samyak-saṃbodhi (unsurpassed perfect awakening).}
\hline

% === LINE 29 ===
\linenum{29} 故知般若波羅蜜多
\newline\rom{gùzhī bōrě bōluómìduō}
&
\deva{तस्माज् ज्ञातव्यम्: प्रज्ञापारमिता}
\newline\rom{tasmāj jñātavyam: prajñāpāramitā}
&
\tib{དེ་ལྟ་བས་ན་ཤེས་རབ་ཀྱི་ཕ་རོལ་ཏུ་ཕྱིན་པའི་སྔགས}
\newline\rom{de lta bas na shes rab kyi pha rol tu phyin pa'i sngags}
\\
\eng{Therefore know that prajñāpāramitā}

% === LINE 30 ===
\linenum{30} 是大神咒，是大明咒
\newline\rom{shì dà shén zhòu, shì dà míng zhòu}
&
\deva{महामन्त्रो महाविद्यामन्त्रो}
\newline\rom{mahā-mantro mahā-vidyā-mantro}
&
\tib{རིག་པ་ཆེན་པོའི་སྔགས}
\newline\rom{rig pa chen po'i sngags}
\\
\eng{is the great mantra, the mantra of great illumination,}

% === LINE 31 ===
\linenum{31} 是無上咒，是無等等咒
\newline\rom{shì wúshàng zhòu, shì wú děngděng zhòu}
&
\deva{ऽनुत्तरमन्त्रो ऽसमसममन्त्रः}
\newline\rom{'nuttara-mantro 'samasama-mantraḥ}
&
\tib{བླ་ན་མེད་པའི་སྔགས། མི་མཉམ་པ་དང་མཉམ་པའི་སྔགས}
\newline\rom{bla na med pa'i sngags. mi mnyam pa dang mnyam pa'i sngags}
\\
\eng{the unsurpassed mantra, the unequalled mantra,}

% === LINE 32 ===
\linenum{32} 能除一切苦，眞實不虛故
\newline\rom{néng chú yīqiè kǔ, zhēnshí bù xū gù}
&
\deva{सर्वदुःखप्रशमनः। सत्यम् अमिथ्यत्वात्}
\newline\rom{sarva-duḥkha-praśamanaḥ. satyam amithyatvāt}
&
\tib{སྡུག་བསྔལ་ཐམས་ཅད་རབ་ཏུ་ཞི་བར་བྱེད་པའི་སྔགས། མི་རྫུན་པས་ན་བདེན་པར་ཤེས་པར་བྱ་སྟེ}
\newline\rom{sdug bsngal thams cad rab tu zhi bar byed pa'i sngags. mi rdzun pas na bden par shes par bya ste}
\\
\eng{which removes all suffering. It is true, not false.}
\hline

% === LINE 33 ===
\linenum{33} 說般若波羅蜜多咒，卽說咒曰
\newline\rom{shuō bōrě bōluómìduō zhòu, jí shuō zhòu yuē}
&
\deva{प्रज्ञापारमितायामुक्तो मन्त्रः। तद्यथा}
\newline\rom{prajñāpāramitāyām ukto mantraḥ. tadyathā}
&
\tib{ཤེས་རབ་ཀྱི་ཕ་རོལ་ཏུ་ཕྱིན་པའི་སྔགས་སྨྲས་པ། ཏདྱ་ཐཱ}
\newline\rom{shes rab kyi pha rol tu phyin pa'i sngags smras pa. tad+ya thA}
\\
\eng{The prajñāpāramitā mantra is proclaimed:}

% === LINE 34: MANTRA ===
\linenum{34} 揭帝揭帝\footnote{\lemma{揭帝} \textit{jiēdì}] Phonetic transcription of \textit{gate}. Mantra transliterated, not translated, in all traditions.}，波羅揭帝，波羅僧揭帝，菩提僧莎訶
\newline\rom{jiēdì jiēdì, bōluó jiēdì, bōluó sēng jiēdì, pútí sēng suōhē}
&
\deva{ॐ गते गते पारगते पारसंगते बोधि स्वाहा}
\newline\rom{Oṃ gate gate pāragate pārasaṃgate bodhi svāhā}
&
\tib{ཨོཾ་ག་ཏེ་ག་ཏེ་པཱ་ར་ག་ཏེ་པཱ་ར་སཾ་ག་ཏེ་བོ་དྷི་སྭཱ་ཧཱ}
\newline\rom{oM ga te ga te pA ra ga te pA ra saM ga te bo d+hi swA hA}
\\
\eng{Gate gate pāragate pārasaṃgate bodhi svāhā!}
\hline

\end{longtable}

\vspace{1em}
\begin{center}
{\small\itshape Heart Sūtra Critical Edition $\cdot$ 34 lines $\cdot$ Chinese Priority (Nattier 1992)}
\end{center}

\end{document}
